% Source: https://artofproblemsolving.com/community/c3h1343062_construction_problem
% Stable identity: aops-topic-1343062-q01
\textbf{Construction.}
\begin{enumerate}
\item Choose any point $A$ on the circle. With the compass set to $r$, mark successively $B$, $C$, $D$ on the circle with $AB=BC=CD=r$.
\item Draw the circle with centre $A$ and radius $AC$, and the circle with centre $D$ and radius $DB$; let $X$ be one of their intersection points.
\item Draw the circle with centre $A$ and radius $OX$. It meets the given circle at $P$ and $Q$. Then $A$, $P$, $D$, $Q$ are the vertices of the square.
\end{enumerate}

\begin{asy}
usepackage("amsmath");
size(5.5cm);
defaultpen(fontsize(10pt));
pair O=(0,0), A=dir(180), B=dir(120), C=dir(60), D=dir(0), X=(0,sqrt(2)), P=dir(90), Q=dir(-90);
draw(unitcircle,linewidth(0.9));
draw(arc(A,sqrt(3),22,aTan(sqrt(2))),blue); draw(arc(D,sqrt(3),180-aTan(sqrt(2)),158),blue);
draw(arc(A,sqrt(2),-55,55),red);
draw(A--P--D--Q--cycle,deepgreen+dashed);
dot(O);dot(A);dot(B);dot(C);dot(D);dot(X);dot(P);dot(Q);
label("$O$",O,S); label("$A$",A,W); label("$B$",B,NW); label("$C$",C,NE); label("$D$",D,E);
label("$X$",X,N); label("$P$",P,NE); label("$Q$",Q,SE);
label("\mbox{}",(min(currentpicture,user=true)+max(currentpicture,user=true))/2,black); // keeps the text after the figure black when LaTeX typesets the labels inline
\end{asy}

\textbf{Proof.} The chords $AB$, $BC$, $CD$ each equal $r$, so each subtends a central angle of $60^\circ$ (triangles $OAB$, $OBC$, $OCD$ are equilateral). Thus $\angle AOD=180^\circ$: $A$ and $D$ are endpoints of a diameter. The chord $AC$ subtends $120^\circ$, so $AC=2r\sin60^\circ=r\sqrt3$, and likewise $DB=r\sqrt3$.

$X$ is equidistant from $A$ and $D$, so $X$ lies on the perpendicular bisector of $AD$, which passes through $O$; hence $\angle XOA=90^\circ$ and, by the Pythagorean theorem,
\[
OX^2=AX^2-OA^2=3r^2-r^2=2r^2,\qquad OX=r\sqrt2 .
\]
Now $AP=r\sqrt2$ with $OA=OP=r$, so $OA^2+OP^2=2r^2=AP^2$ and $\angle AOP=90^\circ$; similarly $\angle AOQ=90^\circ$, with $P$ and $Q$ on opposite sides of $AD$. So $OP$ and $OQ$ are the radii perpendicular to the diameter $AD$, and the points $A,P,D,Q$ divide the circle into four $90^\circ$ arcs: they are the vertices of an inscribed square.
